\subsection{Transforming Domain and Range} 
We can easily find the domain and range of a transformed function. (In fact, we've been doing this all semester!)
\begin{itemize}
\item \makeblank[1.5in] transformations affect the \makeblank[1in]
\item \makeblank[1.5in] transformations affect the \makeblank[1in]
\end{itemize}
\begin{Example}
Below is the graph of $y=f(x)$. Let $g(x)=2 f(3-x)-1$.
\begin{icpic}[x=.5\linespace,y=.5\linespace]
\useasboundingbox (-6,-5) rectangle (6,5);
\setgraph{6}{6}{5}{5}
\GraphSmall
\draw[graphend] (-2,1) -- (1,3) -- (2,-1);
\draw[blue] (1,3) node[anchor=west]{$y=f(x)$};
\end{icpic}
\begin{itemize}
\item $\dom f=\makeblank[1in]$ and $\rng f=\makeblank[1in]$
\end{itemize}
Let's sketch the graph of $g(x)$ to see its domain and range.
\begin{lpic}{}{6}
%\twocolleft{-1}{6}{Method 1}{Method 2}
\end{lpic}
\begin{itemize}
\item $\dom g=\makeblank[1in]$ and $\rng g=\makeblank[1in]$
\item To find the domain, apply only the \makeblank[1.5in] transformations.
\item To find the range, apply only the \makeblank[1.5in] transformations.
\end{itemize}
\begin{lpic}{}{5}
\twocolleft{-1}{5}{$\dom f={}$}{$\rng f={}$}
\twocollblleft{-5}{$\dom g={}$}{$\rng g={}$}
\end{lpic}
\end{Example}